% Sismogramme enregistré par une station : arrivées des ondes P et S (exercice 7, chapitre 17)
\begin{tikzpicture}
\begin{axis}[otMesures,otAxesBout,width=118mm,height=40mm,xmin=0,xmax=45,ymin=-1.3,ymax=1.3,
  axis x line=middle,xtick={0,5,...,45},ytick=\empty,xlabel={$t$ (s)},ylabel={},clip=false]
\addplot[otNavy,line width=0.5pt,domain=0:45,samples=900] {(x>16.8)*0.35*exp(-(x-16.8)/6)*sin(deg(9*(x-16.8))) + (x>28.8)*exp(-(x-28.8)/5)*sin(deg(5*(x-28.8)))};
\solution{
  \draw[otOrange!85!black,line width=0.6pt,dash pattern=on 2pt off 1.5pt] (axis cs:16.8,-1.2) -- (axis cs:16.8,1.0) node[pcLegende,anchor=south,text=otOrange!85!black] {ondes P : \qty{16,8}{s}};
  \draw[otRed!80!black,line width=0.6pt,dash pattern=on 2pt off 1.5pt] (axis cs:28.8,-1.2) -- (axis cs:28.8,1.15) node[pcLegende,anchor=south,text=otRed!80!black] {ondes S : \qty{28,8}{s}};
  \draw[otGreenDark,line width=0.6pt,{Stealth[length=1.4mm]}-{Stealth[length=1.4mm]}] (axis cs:16.8,-1.05) -- (axis cs:28.8,-1.05) node[pos=0.5,pcLegende,anchor=north,text=otGreenDark] {\qty{12}{s}};}
\end{axis}
\end{tikzpicture}
\par\smallskip\legende{Mouvement vertical du sol enregistré par un sismomètre ; l'origine des dates est le début du séisme.\solution{ Les ondes P, plus rapides, arrivent les premières.}}
